\section{Structure-Layer Comparison}
\label{sec:structure}

\subsection{Overview}

The structure layer executes the bisection tree that produces the initial
$k$-district plan.
All chain-based search algorithms start from this initial plan; the initial
plan quality therefore functions as the warm-up baseline.
For SMC and ILP, the structure layer is not used (SMC constructs plans from
scratch; ILP produces a single optimal plan).

We compare six structure algorithms on EC, PP, D-seat, and runtime.
ILP results are reported only for NH ($k = 2$); ILP is not tractable for
larger states.

\subsection{Results}

Table~\ref{tab:structure} gives EC, PP, and runtime for each structure
algorithm on each state.
A dagger (\dag{}) marks new results for FL and NH.

\begin{table}[ht]
\centering
\caption{%
  Structure-layer comparison: edge-cut (\ec{}), mean Polsby-Popper (\pp{}),
  and runtime by algorithm and state.
  EC: lower is more compact. PP: higher is more compact. Runtime in seconds.
  All runs use \texttt{--search convergence} at $T_{\mathrm{conv}} = 600$.
  ILP applies only to NH ($k = 2$).
  Best EC per state is \textbf{bolded}; best PP per state is \textit{italicised}.
  \dag{} denotes new single-run estimates.
}
\label{tab:structure}
\small
\begin{tabular}{@{}lcccccccccccc@{}}
\toprule
 & \multicolumn{3}{c}{\textbf{NC} ($k=14$)}
 & \multicolumn{3}{c}{\textbf{WI} ($k=8$)}
 & \multicolumn{3}{c}{\textbf{TX} ($k=38$)}
 & \multicolumn{3}{c}{\textbf{FL} ($k=28$)\dag{}} \\
\cmidrule(lr){2-4}\cmidrule(lr){5-7}\cmidrule(lr){8-10}\cmidrule(lr){11-13}
Algorithm & EC & PP & t(s) & EC & PP & t(s) & EC & PP & t(s) & EC & PP & t(s) \\
\midrule
METIS          & 4{,}821 & 0.201 &  2.1
               & 2{,}103 & 0.218 &  0.9
               & 14{,}402 & 0.187 &  6.8
               & \dag{}11{,}204 & \dag{}0.179 & \dag{} 4.9 \\
Sim.\ Anneal.  & \textbf{4{,}304} & 0.208 & 25.3
               & \textbf{1{,}879} & 0.224 & 11.1
               & \textbf{12{,}891} & 0.194 & 82.1
               & \dag{}\textbf{10{,}024} & \dag{}0.188 & \dag{}60.4 \\
CVD-Graph      & 4{,}689 & 0.231 &  3.8
               & 2{,}041 & 0.248 &  1.6
               & 14{,}018 & 0.214 & 11.3
               & \dag{}10{,}841 & \dag{}0.209 & \dag{} 8.2 \\
CVD-Geo        & 4{,}632 & \textit{0.249} &  4.2
               & 1{,}998 & \textit{0.267} &  1.8
               & 13{,}741 & \textit{0.232} & 12.7
               & \dag{}10{,}603 & \dag{}\textit{0.228} & \dag{} 9.1 \\
BFS Growth     & 4{,}758 & 0.226 &  1.3
               & 2{,}068 & 0.241 &  0.6
               & 14{,}219 & 0.211 &  4.1
               & \dag{}11{,}003 & \dag{}0.203 & \dag{} 2.9 \\
\bottomrule
\end{tabular}
\end{table}

\begin{table}[ht]
\centering
\caption{%
  Structure-layer results for New Hampshire ($k=2$, $N \approx 260$)\dag{}.
  ILP achieves certified-optimal EC; all other algorithms are heuristics.
  Best EC is \textbf{bolded}; best PP is \textit{italicised}.
}
\label{tab:structure-nh}
\small
\begin{tabular}{@{}lccc@{}}
\toprule
Algorithm & EC (NH) & PP (NH) & Runtime (s) \\
\midrule
METIS              & 58 & 0.341 &  0.2 \\
Sim.\ Anneal.      & 51 & 0.348 &  1.9 \\
CVD-Graph          & 54 & 0.362 &  0.3 \\
CVD-Geo            & 53 & \textit{0.374} &  0.3 \\
BFS Growth         & 56 & 0.358 &  0.1 \\
ILP (exact)        & \textbf{47} & 0.371 & 12.8 \\
\bottomrule
\end{tabular}
\end{table}

\subsection{ILP: Certified Optimal for NH}

\ilp{} (B.24) achieves \ec{} = 47 for NH, certified as the global minimum
by the ILP optimality certificate (branch-and-bound concluded with zero
integrality gap).
This is 8.5\% below Simulated Annealing (51), the nearest heuristic, and
23.4\% below METIS (58).

The ILP optimality gap across heuristics reveals that even the best heuristic
structure algorithm leaves meaningful compactness on the table for small
instances.
For larger states (NC, WI, TX, FL), ILP is not applicable due to the
exponential growth of the branch-and-bound tree with $N$.
ILP's \pp{} (0.371) is competitive with CVD-Geo (0.374) for NH;
\pp{} maximisation is not an ILP objective, so this reflects the incidental
geometric regularity of minimum-\ec{} bisections on small graphs.

\subsection{Simulated Annealing: Best Heuristic EC}

\sa{} (B.19) achieves the lowest \ec{} among heuristic structure algorithms
in all states: 10.7\% below METIS for NC, 10.6\% for WI, 10.5\% for TX,
and 10.5\% for FL.
The improvement is consistent across state sizes, reflecting that SA's
node-swap improvement steps operate at a fixed scale independent of $k$.

The cost is runtime: SA is 12$\times$ slower than METIS for NC (25.3 s
vs.\ 2.1 s), and proportionally slower for larger states (82.1 s for TX
vs.\ 6.8 s for METIS).
For a 50-state sweep, this translates to approximately 10$\times$ the
structure-layer compute budget.
When \ec{} optimisation is the primary goal and runtime is not constrained,
SA is the correct structure choice.

SA's \pp{} is only marginally better than METIS (0.208 vs.\ 0.201 for NC):
SA optimises \ec{} directly; \pp{} improvement is a side effect of fewer
cross-district adjacencies, not an explicit objective.

\subsection{CVD-Geographic: Best Heuristic PP}

\cvdgeo{} (B.22b) achieves the highest mean \pp{} among all heuristic
structure algorithms across all states: 0.249 for NC, 0.267 for WI, 0.232
for TX, 0.228 for FL.
This exceeds even the best search-algorithm results for \pp{} (Section
\ref{sec:search}), confirming that the structure layer sets the \pp{} ceiling.

\cvdgeo{} requires tract centroid data (latitude/longitude of the population-
weighted centroid of each census tract).
For the 2020 decennial census, centroid data is available in the
\texttt{INTPTLAT}/\texttt{INTPTLON} fields of the census tract shapefile;
the \texttt{redist} data pipeline populates these fields automatically during
the \texttt{redist fetch} step.
No additional data acquisition is required for users running the standard
pipeline.

\cvdgeo{}'s \ec{} (4{,}632 for NC) is 3.9\% below METIS but 7.4\% above SA.
This is the fundamental CVD-Geo trade-off: superior geometric compactness
(\pp{}) at the cost of suboptimal graph compactness (\ec{}).
For tasks where the \pp{} criterion is primary (as in many legal standards
that reference ``compactness'' in the geometric rather than graph-theoretic
sense), CVD-Geo is the correct structure choice.

\subsection{CVD Graph-Distance: PP Without Centroids}

\cvdg{} (B.22) achieves \pp{} = 0.231 for NC, below CVD-Geo (0.249) but
substantially above METIS (0.201).
The graph-hop metric approximates Euclidean distance on planar graphs;
the approximation degrades for states with highly irregular tract shapes
(coastal FL: 0.209 for CVD-Graph vs.\ 0.228 for CVD-Geo).
CVD-Graph requires only the adjacency graph and is therefore applicable
without centroid data, making it the preferred choice for custom graph
inputs or non-standard geographic units.

\subsection{BFS Growth: Fastest Structure Method}

\bfsgrow{} (B.23) runs in wall-clock time 1.3 s for NC, 0.6 s for WI,
4.1 s for TX, 2.9 s for FL --- the fastest structure method by 1.6--2$\times$
over METIS.
\pp{} values (0.226 for NC, 0.241 for WI, 0.211 for TX, 0.203 for FL) are
substantially above METIS and close to CVD-Graph, reflecting that BFS growth
from random seeds produces geographically contiguous districts with natural
boundaries.

BFS Growth has zero hyperparameters: the only choice is the number of random
seeds, which equals $k$.
This makes it reproducible and robust across states.
\ec{} performance is near METIS (within 1\%); BFS Growth does not explicitly
optimise \ec{}.
For applications where structure runtime is a bottleneck and moderate \pp{}
is required without the centroid data of CVD-Geo, BFS Growth is the
appropriate choice.

\subsection{Summary: Structure Layer}

The structure-layer hierarchy, consistent across all states:

\paragraph{By EC:}
\[
  \ilp{} < \sa{} < \cvdgeo{} < \cvdg{} < \bfsgrow{} \lesssim \metis{}
\]
where $<$ denotes a clear gap and $\lesssim$ denotes within 1\%.

\paragraph{By PP:}
\[
  \cvdgeo{} > \cvdg{} > \bfsgrow{} \approx \sa{} > \metis{}
\]

\paragraph{By runtime:}
\[
  \bfsgrow{} < \metis{} < \cvdg{} \approx \cvdgeo{} \ll \sa{} \ll \ilp{}
\]

No single algorithm dominates all three dimensions.
The decision framework in Section~\ref{sec:decision} uses these three
orderings to route practitioners to the correct structure algorithm
based on their primary criterion.
